
\documentclass{article}
\usepackage{colortbl}
\usepackage{makecell}
\usepackage{multirow}
\usepackage{supertabular}

\begin{document}

\newcounter{utterance}

\centering \large Interaction Transcript for game `eqbench', experiment `eq\_bench', episode 0 with reflexion{-}Qwen3.5{-}9B.
\vspace{24pt}

{ \footnotesize  \setcounter{utterance}{1}
\setlength{\tabcolsep}{0pt}
\begin{supertabular}{c@{$\;$}|p{.15\linewidth}@{}p{.15\linewidth}p{.15\linewidth}p{.15\linewidth}p{.15\linewidth}p{.15\linewidth}}
    \# & \multicolumn{2}{c}{Player} && \multicolumn{2}{c}{Game Master} \\
    \hline

    \theutterance \stepcounter{utterance}  
    & & & \multicolumn{4}{p{0.6\linewidth}}{
        \cellcolor[rgb]{0.9,0.9,0.9}{
            \makecell[{{p{\linewidth}}}]{
                \texttt{\tiny{[P1$\langle$GM]}}
                \\ 
\texttt{Your task is to predict the likely emotional responses of a character in this dialogue:} \\
\\ 
\texttt{Robert: Claudia, you've always been the idealist. But let's be practical for once, shall we?} \\
\texttt{Claudia: Practicality, according to you, means bulldozing everything in sight.} \\
\texttt{Robert: It's called progress, Claudia. It's how the world works.} \\
\texttt{Claudia: Not my world, Robert.} \\
\texttt{Robert: Your world? You mean this...this sanctuary of yours?} \\
\texttt{Claudia: It's more than a sanctuary. It's a testament to our parents' love for nature.} \\
\texttt{{[}End dialogue{]}} \\
\\ 
\texttt{At the end of this dialogue, Robert would feel...} \\
\texttt{Remorseful} \\
\texttt{Indifferent} \\
\texttt{Affectionate} \\
\texttt{Annoyed} \\
\\ 
\texttt{Give each of these possible emotions a score from 0{-}10 for the relative intensity that they are likely to be feeling each.} \\
\\ 
\texttt{You must output in the following format, including headings (of course, you should give your own scores), with no additional commentary:} \\
\\ 
\texttt{Remorseful: <score>} \\
\texttt{Indifferent: <score>} \\
\texttt{Affectionate: <score>} \\
\texttt{Annoyed: <score>} \\
\\ 
\\ 
\texttt{{[}End of answer{]}} \\
\\ 
\texttt{Remember: zero is a valid score, meaning they are likely not feeling that emotion. You must score at least one emotion > 0.} \\
\\ 
\texttt{Your answer:} \\
            }
        }
    }
     \\ \\

    \theutterance \stepcounter{utterance}  
    & \multicolumn{4}{p{0.6\linewidth}}{
        \cellcolor[rgb]{0.561,0.737,0.561}{
            \makecell[{{p{\linewidth}}}]{
                \texttt{\tiny{[P1$\rangle$GM]}}
                \texttt{Remorseful: 2} \\
\texttt{Indifferent: 1} \\
\texttt{Affectionate: 0} \\
\texttt{Annoyed: 9} \\
            }
        }
    }
    & & \\ \\

    \theutterance \stepcounter{utterance}  
    & & & \multicolumn{2}{p{0.3\linewidth}}{
        \cellcolor[rgb]{0.9,0.9,0.9}{
            \makecell[{{p{\linewidth}}}]{
                \texttt{\tiny{[GM$|$GM]}}
                \texttt{Remorseful: 2} \\
\texttt{Indifferent: 1} \\
\texttt{Affectionate: 0} \\
\texttt{Annoyed: 9} \\
            }
        }
    }
    & & \\ \\

    \theutterance \stepcounter{utterance}  
    & & & \multicolumn{2}{p{0.3\linewidth}}{
        \cellcolor[rgb]{0.9,0.9,0.9}{
            \makecell[{{p{\linewidth}}}]{
                \texttt{\tiny{[GM$|$GM]}}
                \texttt{Remorseful: 0} \\
\texttt{Indifferent: 6} \\
\texttt{Affectionate: 0} \\
\texttt{Annoyed: 7} \\
            }
        }
    }
    & & \\ \\

    \theutterance \stepcounter{utterance}  
    & & & \multicolumn{2}{p{0.3\linewidth}}{
        \cellcolor[rgb]{0.9,0.9,0.9}{
            \makecell[{{p{\linewidth}}}]{
                \texttt{\tiny{[GM$|$GM]}}
                \texttt{game\_result = LOSE} \\
            }
        }
    }
    & & \\ \\

    \theutterance \stepcounter{utterance}  
    & & \multicolumn{4}{p{0.6\linewidth}}{
        \cellcolor[rgb]{0.95,0.95,0.95}{
            \makecell[{{p{\linewidth}}}]{
                \texttt{\tiny{[P1$\leftrightarrow$P]}}
                \texttt{Post{-}episode Reflexion (not sent to GM):} \\
\texttt{New plan: I incorrectly scored "Annoyed" as 9 and "Remorseful" as 2, failing to recognize that Robert's defensive tone ("bulldozing everything," "progress") stems from frustration and irritation rather than genuine regret or indifference. His focus on "practicality" vs. "idealist" labels indicates a conflict where he feels his pragmatic approach is being rejected, leading to high annoyance. Conversely, there is no evidence of remorse (he doesn't apologize) or affection (the tone is combative). Next time, I will analyze the specific vocabulary: "bulldozing" and "progress" signal aggression and frustration, warranting a high score for Annoyed (9{-}10) and low scores for Remorseful (0{-}1) and Affectionate (0). I will ensure the sum of scores reflects the dominant negative emotion without artificially inflating unrelated traits.} \\
\\ 
\texttt{Remorseful: 1} \\
\texttt{Indifferent: 0} \\
\texttt{Affectionate: 0} \\
\texttt{Annoyed: 9} \\
            }
        }
    }
    & \\ \\

\end{supertabular}
}

\end{document}
